% Generated draft: requires mathematical and source review.
\documentclass[11pt]{article}
\usepackage[margin=25mm]{geometry}
\usepackage{fontspec}
\setmainfont{Latin Modern Roman}
\usepackage{amsmath,amssymb,mathtools,unicode-math}
\setmathfont{Latin Modern Math}
\usepackage{enumitem,needspace,xurl,hyperref}
\setlength{\parindent}{0pt}
\setlength{\parskip}{5pt}
\setlength{\emergencystretch}{3em}
\newcommand{\N}{\mathbb{N}}
\newcommand{\Z}{\mathbb{Z}}
\newcommand{\Q}{\mathbb{Q}}
\newcommand{\R}{\mathbb{R}}
\newcommand{\C}{\mathbb{C}}
\newcommand{\F}{\mathbb{F}}
\newcommand{\A}{\mathbb{A}}
\newcommand{\PP}{\mathbb{P}}
\newcommand{\E}{\mathbb{E}}
\newcommand{\eps}{\varepsilon}
\newcommand{\dd}{\,\mathrm{d}}
\begin{document}
\section*{UnsolvedMath ID: 801154}
\textbf{Source problem code:} catalogue entry
\subsection{Definitions and mathematical statement}
Under the conventional real-manifold interpretation, let \(M\) be a two-dimensional manifold and let \(L\colon TM\to TM\) be a field of endomorphisms. The supplied excerpt does not specify the differentiability category. Its Nijenhuis torsion is
\[
N_L(X,Y)=[LX,LY]-L[LX,Y]-L[X,LY]+L^2[X,Y]
\]
for vector fields \(X,Y\). The operator \(L\) is Nijenhuis if \(N_L(X,Y)=0\) for all \(X,Y\). It is \(\mathrm{gl}\)-regular if, at every \(p\in M\), the minimal polynomial of \(L_p\) has degree two; equivalently, \(L_p\) is not a scalar multiple of the identity.

For precision, the conventional differential interpretation of geodesic compatibility is recorded here. A nondegenerate symmetric metric tensor \(g\), with Levi--Civita connection \(\nabla\), is geodesically compatible with \(L\) if
\[
g(LX,Y)=g(X,LY)
\]
and
\[
(\nabla_XL)Y=\frac12\left(\dd(\operatorname{tr}L)(Y)X+g(X,Y)\operatorname{grad}_g(\operatorname{tr}L)\right)
\]
for all vector fields \(X,Y\), where \(g(\operatorname{grad}_g h,Z)=\dd h(Z)\). This convention permits indefinite metrics; the excerpt does not specify the intended signature or provide its own definition of compatibility.

\textbf{Requested construction.} For every \(\mathrm{gl}\)-regular Nijenhuis operator \(L\) in dimension two, construct explicitly a metric \(g\) geodesically compatible with \(L\).

The supplied text suggests first obtaining a regular conservation law \(f\), then constructing companion coordinates and applying an explicit formula cited as \texttt{[nijappl5]}. This is a suggested route, not an additional hypothesis on \(L\). The excerpt supplies neither the definitions needed for that route nor the cited formula. It also does not determine whether the requested construction is local near every point or global on \(M\), the permitted metric signatures, or the regularity category. These qualifications remain unresolved; the displayed compatibility equations express a conventional interpretation requiring confirmation against the complete source context.
\subsection{Short English statement}
For each two-dimensional \(\mathrm{gl}\)-regular Nijenhuis operator, explicitly construct a geodesically compatible metric. The source suggests using a regular conservation law and companion coordinates. Local versus global scope, metric signature, regularity, and the precise source convention for compatibility require clarification.
\subsection{Sources}
\begin{itemize}
\item Local UnsolvedMath registry record: permanent ID \(801154\), exact source problem code \texttt{catalogue entry}, title ``Nijenhuis Geometry and Integrable Systems --- Problem 7.'' The supplied registry labels the problem open; this label is not an independently verified mathematical or literature assessment.
\item Bolsinov, Konyaev, and Matveev, \emph{Research Problems on Nijenhuis Geometry and Integrable Systems} (2024), Problem 7. Supplied source URL: \url{https://arxiv.org/abs/2410.04276}. Attribution and proposal year are supplied by the record.
\item Supplied provenance records access on 2026-10-16 and extraction by \texttt{arxiv\_latex}. These are supplied provenance statements, not evidence of independent source consultation for this draft.
\item The statement cites an explicit formula using the reference label \texttt{[nijappl5]}. Its bibliographic entry and formula are absent from the supplied record, so this reference cannot be identified more precisely here.
\end{itemize}
\end{document}
